\documentclass[10pt,a4paper]{amsart}
\usepackage[margin=19mm]{geometry}
\usepackage{amsmath,amssymb,mathtools}
\usepackage[scr=rsfs,cal=euler]{mathalfa}
\usepackage{xfrac}
\usepackage[unicode,hidelinks,bookmarks=false]{hyperref}
\newcommand{\ie}{i.e.\ }
\newcommand{\cf}{cf.\ }
\newcommand{\resp}{respectively\ }
\newcommand{\Subsection}{\subsection}
\providecommand{\nobreakdash}{\nobreak\hbox{-}}
\setlength{\emergencystretch}{2em}
\multlinegap0pt
\usepackage{float}
\usepackage{xcolor}
\usepackage{amssymb}
\usepackage{mathtools}
\usepackage[shortlabels]{enumitem}
\newtheorem{theorem}{Theorem}
\title{SMI: Statistical Membership Inference for Reliable Unlearned Model Auditing}
\author{Jialong Sun, Zeming Wei, Jiaxuan Zou, Jiacheng Gong, Jie Fu, Chengyang Dong, Heng Xu, Jialong Li, Bo Liu}
\date{Source: arXiv:2602.01150v2 (CC BY 4.0)}
\begin{document}
\maketitle

\noindent\emph{Source and licence.} SMI: Statistical Membership Inference for Reliable Unlearned Model Auditing. Version arXiv:2602.01150v2 (CC BY 4.0), distributed by arXiv under the Creative Commons Attribution 4.0 International licence, http://creativecommons.org/licenses/by/4.0/, as recorded in the arXiv licence metadata for this exact version and shown on the abstract page https://arxiv.org/abs/2602.01150v2. Full licence notice: \texttt{licenses/arxiv-2602.01150-CC-BY-4.0.txt}.\par\smallskip

\noindent\textit{Editorial scope.} Counted unit: the article's theorem theorem (source0.tex lines 736--745). The article's complete original proof of it is delivered with it (source0.tex lines 747--923, one \texttt{proof} environment of 177 lines). No mathematics is written, repaired or re-derived: the statement and the proof are the article's own lines, in the article's own notation, and no retained line is substituted or re-flowed.  No further source-local context is needed: every cross-reference of every retained line resolves inside the two delivered spans.  Nothing was omitted from any retained span, so the package carries no omission record.  No retained line contains a citation, so the wrapper carries no bibliography and no bibliography entry.  No retained line contains a figure environment, an image-inclusion command or drawing code, so no image is delivered and the package declares no asset dependency.  Editorial formatting only, in addition: an AMS \texttt{amsart} preamble that restates the article's own declarations and macros used by the retained lines, namely the environments \texttt{theorem} on separate counters, together with the standard packages those lines need.  The retained environments print in this document as Theorem 1 in order of appearance, whereas the article prints the same items with the numbering of its own full document.  The complete native source of the article as distributed in the arXiv e-print for this exact version is bundled unchanged as \texttt{sources/193763/source0.tex}.
\begin{theorem}[PAC-Bayesian Generalization Bound with $\chi^2$ Divergence]
Let $\mathcal{H}$ be a hypothesis class and $\ell: \mathcal{H} \times \mathcal{Z} \to [0,1]$ a loss function. Let $D$ be a data distribution and $S = (z_1, \dots, z_m) \sim D^m$ an i.i.d. sample. For any prior distribution $\mathcal{P}$ over $\mathcal{H}$ and any posterior distribution $\mathcal{Q}$ such that $\mathcal{Q} \ll \mathcal{P}$, with probability at least $1 - \delta$ over the draw of $S$, it holds that
\begin{equation}
    R_D(\mathcal{Q}) \leq R_S(\mathcal{Q}) + \sqrt{ \frac{2}{m\delta} \left( \chi^2(\mathcal{Q} \| \mathcal{P}) + 1 \right)  },
\end{equation}
where $R_D(\mathcal{Q}) = \mathbb{E}_{h \sim \mathcal{Q}}[R_D(h)]$, $R_S(\mathcal{Q}) = \mathbb{E}_{h \sim \mathcal{Q}}[R_S(h)]$, and the $\chi^2$-divergence is defined as
\begin{equation}
    \chi^2(\mathcal{Q} \| \mathcal{P}) = \int_{\mathcal{H}} \left( \frac{d\mathcal{Q}(h)}{d\mathcal{P}(h)} - 1 \right)^2 d\mathcal{P}(h).
\end{equation}
\end{theorem}

\begin{proof}
Let $S=(z_1,\ldots,z_m)\sim D^m$ be a training sample drawn independently from the distribution $D$. For any hypothesis $h\in\mathcal H$, define its true risk and empirical risk as
\begin{equation}
R_D(h)=\mathbb E_{z\sim D}[\ell(h,z)],
\qquad
R_S(h)=\frac{1}{m}\sum_{i=1}^m \ell(h,z_i).
\end{equation}
We further define the generalization gap of a single hypothesis as
\begin{equation}
\Delta_S(h)=R_D(h)-R_S(h).
\end{equation}
Since
\begin{equation}
\mathbb E_{S\sim D^m}[R_S(h)]
=
R_D(h),
\end{equation}
we have, for any fixed $h\in\mathcal H$,
\begin{equation}
\mathbb E_{S\sim D^m}[\Delta_S(h)]=0.
\end{equation}

For a posterior distribution $\mathcal Q$, the corresponding Gibbs true risk and empirical risk are given by
\begin{equation}
R_D(\mathcal Q)
=
\mathbb E_{h\sim\mathcal Q}[R_D(h)],
\qquad
R_S(\mathcal Q)
=
\mathbb E_{h\sim\mathcal Q}[R_S(h)].
\end{equation}
Therefore,
\begin{equation}
R_D(\mathcal Q)-R_S(\mathcal Q)
=
\mathbb E_{h\sim\mathcal Q}[\Delta_S(h)].
\end{equation}

Assume that $\mathcal Q\ll \mathcal P$, and let
\begin{equation}
r(h)=\frac{d\mathcal Q}{d\mathcal P}(h)
\end{equation}
denote the Radon--Nikodym derivative of $\mathcal Q$ with respect to $\mathcal P$. Then,
\begin{equation}
\mathbb E_{h\sim\mathcal Q}[\Delta_S(h)]
=
\int_{\mathcal H}\Delta_S(h)r(h)\,d\mathcal P(h)
=
\mathbb E_{h\sim\mathcal P}[r(h)\Delta_S(h)].
\end{equation}
By the Cauchy--Schwarz inequality, we obtain
\begin{equation}
\mathbb E_{h\sim\mathcal P}[r(h)\Delta_S(h)]
\le
\left|
\mathbb E_{h\sim\mathcal P}[r(h)\Delta_S(h)]
\right|
\le
\sqrt{\mathbb E_{h\sim\mathcal P}[r(h)^2]}
\sqrt{\mathbb E_{h\sim\mathcal P}[\Delta_S(h)^2]}.
\end{equation}
By the definition of the $\chi^2$-divergence,
\begin{equation}
\chi^2(\mathcal Q\|\mathcal P)
=
\int_{\mathcal H}(r(h)-1)^2\,d\mathcal P(h)
=
\mathbb E_{h\sim\mathcal P}[r(h)^2]-1,
\end{equation}
where we used
\begin{equation}
\mathbb E_{h\sim\mathcal P}[r(h)]=1.
\end{equation}
Hence,
\begin{equation}
\mathbb E_{h\sim\mathcal P}[r(h)^2]
=
\chi^2(\mathcal Q\|\mathcal P)+1.
\end{equation}
Let
\begin{equation}
Z_S
=
\mathbb E_{h\sim\mathcal P}[\Delta_S(h)^2].
\end{equation}
Then,
\begin{equation}
R_D(\mathcal Q)-R_S(\mathcal Q)
\le
\sqrt{\chi^2(\mathcal Q\|\mathcal P)+1}
\sqrt{Z_S}.
\end{equation}

It remains to derive a high-probability upper bound for $Z_S$. Since the prior $\mathcal P$ is independent of the sample $S$, by Fubini's theorem we have
\begin{equation}
\mathbb E_{S\sim D^m}[Z_S]
=
\mathbb E_{h\sim\mathcal P}
\mathbb E_{S\sim D^m}[\Delta_S(h)^2].
\end{equation}
For any fixed $h$, since $\mathbb E_S[\Delta_S(h)]=0$, we have
\begin{equation}
\mathbb E_{S\sim D^m}[\Delta_S(h)^2]
=
\operatorname{Var}_{S\sim D^m}(R_S(h)).
\end{equation}
Moreover,
\begin{equation}
R_S(h)=\frac{1}{m}\sum_{i=1}^m \ell(h,z_i).
\end{equation}
Since $z_1,\ldots,z_m$ are independent and identically distributed,
\begin{equation}
\operatorname{Var}_{S\sim D^m}(R_S(h))
=
\frac{1}{m^2}
\sum_{i=1}^m
\operatorname{Var}_{z_i\sim D}(\ell(h,z_i)).
\end{equation}
By the bounded-variance assumption, there exists a constant $\sigma^2>0$ such that, for all $h\in\mathcal H$,
\begin{equation}
\operatorname{Var}_{z\sim D}(\ell(h,z))\le \sigma^2.
\end{equation}
Therefore,
\begin{equation}
\operatorname{Var}_{S\sim D^m}(R_S(h))
\le
\frac{\sigma^2}{m}.
\end{equation}
Consequently,
\begin{equation}
\mathbb E_{S\sim D^m}[Z_S]
\le
\frac{\sigma^2}{m}.
\end{equation}

Since $Z_S\ge 0$, Markov's inequality gives
\begin{equation}
\Pr_{S\sim D^m}
\left(
Z_S
\ge
\frac{\sigma^2}{m\delta}
\right)
\le
\frac{\mathbb E_S[Z_S]}{\sigma^2/(m\delta)}
\le
\delta.
\end{equation}
Thus, with probability at least $1-\delta$,
\begin{equation}
Z_S
\le
\frac{\sigma^2}{m\delta}.
\end{equation}
Substituting this bound into the previous Cauchy--Schwarz inequality yields that, with probability at least $1-\delta$, for all $\mathcal Q\ll\mathcal P$,
\begin{equation}
R_D(\mathcal Q)-R_S(\mathcal Q)
\le
\sqrt{\chi^2(\mathcal Q\|\mathcal P)+1}
\sqrt{\frac{\sigma^2}{m\delta}}.
\end{equation}
Equivalently,
\begin{equation}
R_D(\mathcal Q)
\le
R_S(\mathcal Q)
+
\sqrt{
\frac{\sigma^2}{m\delta}
\left(
\chi^2(\mathcal Q\|\mathcal P)+1
\right)
}.
\end{equation}
This completes the proof.
\end{proof}

\end{document}
